% HOW A \relax STANDS IN FOR A CONDITIONAL.
%
% An \else, \or or \fi met while something is scanning a number stands for
% \relax: it terminates the scan and stays in the input for the conditional
% to find afterwards. The reference backs the two up ONE AT A TIME -- the
% token first, the \relax over it -- so the frame the \relax stood in is
% spent and gone by the time a fault below it is shown:
%
%   ! Missing number, treated as zero.
%   <to be read again>
%                      \relax
%   <to be read again>
%                      \fi
%
% This engine pushed the pair as one inserted text, so the second frame
% still held the \fi and drew `<inserted text> \relax' with the \fi behind
% the reading.
%
% trip line 419 writes \ifcase1\or\ifeof\fi.
\catcode`\{=1 \catcode`\}=2
\tracingonline=1 \errorcontextlines=5
\message{[A a \fi where a number should be]}
\ifcase1 \or\ifeof\fi\fi\fi
\message{[B an \else where a number should be]}
\ifcase1 \or\ifnum\else\fi\fi\fi
\message{[done]}
\end
